\section{Order of Group and its elements}
Order of a group $G$, denoted as $|G|$ is the number of its element. If $|G|$ is a finite, then $G$ is finite group; otherwise it is infinte group which has infinte order. $(R, +)$ has infinte order; whereas symmetries of a triangle , $\Z/n\Z$ or $(\Z/n\Z)^{\times}$ has finite order.
\begin{note}
$\Z/n\Z$ is a group under addition and $(\Z/n\Z)^{\times}$ is a group under multiplication, consists of the elements of $\Z/n\Z$ that do have multiplicative inverse. Former is called \textbf{additive group} and the latter is called \textbf{multiplicative group}.
\end{note}
\begin{eg}
    $|(\znz)^\times| = \{a \in \znz : \text{gcd}(a, n) = 1\}  =\varphi(n).$\\
    The function is called \textbf{Euler's totient function}.
\end{eg}
\begin{definition}
    For a group $G$ and $x \in G$, the \textit{order} of $x$ is the smallest positive integer $n$ such that $x^n = e$, and is denoted by $|x|$. In this case, $x$ is said to be order of $n$. If no positive power $x$ is the identity, the order of $x$ is defined to be infinity and $x$ is said to be of infinite order.
\end{definition}
\begin{eg}
    $\znz$ has order $6$. Elements of $\znz$ has following orders: $|\overline{0}|=1, |\overline{1}|=6, |\overline{2}|=3, |\overline{3}|=2, |\overline{4}|=3, |\overline{5}|=6$. Notice, order of any element is less than or equals to the order of the group.
\end{eg}
\begin{eg}
    In additive groups $\Z, \Q, \R$ or $\C$, every nonzero element has infinte order. whereas, in the multiplicative groups $\R^\times$ or $\Q^\times$ the element $-1$ has order $2$ and all other nonidentity elements have infinite order.
\end{eg}
\begin{eg}
    In the group $GL_n(R)$, order of
    $\begin{pmatrix}
    -1 & 0\\
    0 & -1
    \end{pmatrix}$ is $2$ because
    $\begin{pmatrix}
    -1 & 0\\
    0 & -1
    \end{pmatrix}^2
    = \begin{pmatrix}
    1 & 0\\
    0 & 1
    \end{pmatrix}$.
\end{eg}
\section{Cyclic Groups}
\begin{definition}
    A group $H$ is cyclic if $H$ can be generated by a single element, i.e., there is some element $x \in H$ such that $H = \{x^n : n \in \Z\}$ (where as usual the operation is multiplication).  
\end{definition}
In additive notion $H$ is cyclic if $H = \{nx : n \in \Z\}$. In both cases we will write $H = \langle x \rangle$ and say $H$ is generated by $x$.
\begin{eg}[A cyclic group may have more than one generator]
    \begin{itemize}
        \item $\znz = \langle \overline{1} \rangle = \langle \overline{2} \rangle$
        \item $\Z/ p \Z = \langle \overline{1} \rangle = ... = \langle \overline{p-1} \rangle$
        \item $(\Z, +) = \langle \overline{-1}\rangle = \langle \overline{1}\rangle$ 
    \end{itemize}
\end{eg}
\begin{eg}
    The 6th complex roots of unity form a cyclic group under multiplication. Here, $z$ is a generator, but $z^2$ is not, because its power fail to produce the odd powers of $z$.   
    \begin{figure}[H]
        \centering
        \incfig{6th-root-of-unity}
        \caption{The 6th complex roots of unity form a cylic group under multiplication}
        \label{fig:6th-root-of-unity}
    \end{figure}
\end{eg}
\section{Centralizers and Normalizers}
\begin{definition}[Centralizer]
    The \textit{centralizer} of a subset $A$ of group $G$ is the set $C_G(A) = \{g \in G : g a g^{-1} = a \text{ for all } a \in A\}$. Since $g a g^{-1} = a$ if and only if $ga = ag, C_G(A)$ is the set of elements of $G$ which commute with every element of $A$.    
\end{definition}
We show $C_G(A)$ is a subgroup of $G$. First of all, $C_G(A) \neq \emptyset$ as $e \in C_G(A)$ because the identity commutes with every element of $G$, in particular for all $a \in A$.\\
Let $x, y \in C_G(A)$, that is, for all $a \in A$
$$
x a x^{-1} = a
\text{ and } 
y a y^{-1} = a.
$$
We have to prove $xy^{-1}$ also exits in $C_G(A)$.
\begin{align*}
    (xy^{-1})  a (xy^{-1})^{-1} &= (xy^{-1})  a (y x^-1)\\
                            &= x (y a y ^{-1}) x^{-1}\\
                            &= x a x^{-1}\\
                            &=a
\end{align*}
In special case when $A = {a}$ we will write simply $C_G(a)$ instead of $C_G({a})$. In this case $a^n \in C_G(a)$ for all $n \in \Z$.\\   
\textbf{To be continued...}